\documentclass[10pt,a4paper]{amsart}
\usepackage[margin=19mm]{geometry}
\usepackage{amsmath,amssymb,mathtools}
\usepackage[scr=rsfs,cal=euler]{mathalfa}
\usepackage{xfrac}
\usepackage[unicode,hidelinks,bookmarks=false]{hyperref}
\newcommand{\ie}{i.e.\ }
\newcommand{\cf}{cf.\ }
\newcommand{\resp}{respectively\ }
\newcommand{\Subsection}{\subsection}
\providecommand{\nobreakdash}{\nobreak\hbox{-}}
\setlength{\emergencystretch}{2em}
\multlinegap0pt
\usepackage{tikz}
\usetikzlibrary{cd}
\usepackage{amssymb}
\usepackage{bm}
\usepackage{tikz}
\newtheorem{lemma}{Lemma}
\title{\textbf{Aspherical Lagrangian submanifolds, Audin's conjecture and cyclic dilations}}
\author{Yin Li}
\date{Source: arXiv:2308.05086v6 (CC BY 4.0)}
\begin{document}
\maketitle

\noindent\emph{Source and licence.} \textbf{Aspherical Lagrangian submanifolds, Audin's conjecture and cyclic dilations}. Version arXiv:2308.05086v6 (CC BY 4.0), distributed by arXiv under the Creative Commons Attribution 4.0 International licence, http://creativecommons.org/licenses/by/4.0/, as recorded in the arXiv licence metadata for this exact version and shown on the abstract page https://arxiv.org/abs/2308.05086v6. Full licence notice: \texttt{licenses/arxiv-2308.05086-CC-BY-4.0.txt}.\par\smallskip

\noindent\textit{Editorial scope.} Counted unit: the article's lemma lemma:DK (source0.tex lines 625--631). The article's complete original proof of it is delivered with it (source0.tex lines 632--663, one \texttt{proof} environment of 32 lines). No mathematics is written, repaired or re-derived: the statement and the proof are the article's own lines, in the article's own notation, and no retained line is substituted or re-flowed.  Retained uncounted as the source-local context the proof needs: lines 594--596; lines 597--599; lines 600--604.  Nothing was omitted from any retained span, so the package carries no omission record.  No retained line contains a citation, so the wrapper carries no bibliography and no bibliography entry.  The retained lines contain the article's own diagram code, which the wrapper typesets with the standard tikz packages; no image file is delivered and the package declares no asset dependency.  Editorial formatting only, in addition: an AMS \texttt{amsart} preamble that restates the article's own declarations and macros used by the retained lines, namely the environments \texttt{lemma} on separate counters, together with the standard packages those lines need.  The retained environments print in this document as Lemma 1 in order of appearance, whereas the article prints the same items with the numbering of its own full document.  The complete native source of the article as distributed in the arXiv e-print for this exact version is bundled unchanged as \texttt{sources/145431/source0.tex}.
\begin{equation}\label{eq:cod}
\delta_{k+1,j}\circ\delta_{k,i}=\delta_{k+1,i}\circ\delta_{k,j-1}\textrm{ for }i<j,
\end{equation}

\begin{equation}\label{eq:cos}
\sigma_{k-1,j}\circ\sigma_{k,i}=\sigma_{k-1,i}\circ\sigma_{k,j+1}\textrm{ for }i\leq j,
\end{equation}

\begin{equation}\label{eq:sd}
\sigma_{k,j}\circ\delta_{k+1,i}=\left\{\begin{array}{ll}
\delta_{k,i}\circ\sigma_{k-1,j-1} & i<j, \\ \mathit{id}_C & i=j,j+1, \\ \delta_{k,i-1}\circ\sigma_{k-1,j} & i>j+1.
\end{array}\right.
\end{equation}

\begin{lemma}\label{lemma:DK}
The composition
\begin{equation}
\psi_{DK}:C_\ast^\mathit{nm}\hookrightarrow C_\ast\twoheadrightarrow C_\ast^\mathit{nd}
\end{equation}
is an isomorphism of chain complexes. In particular, the projection map $C_\ast\twoheadrightarrow C_\ast^\mathit{nd}$ is also a quasi-isomorphism.
\end{lemma}

\begin{proof}
It suffices to show that the natural map $C_\ast^\mathit{nm}(k)\rightarrow C_\ast^\mathit{nd}(k)$ is an isomorphism for each $k\in\mathbb{Z}_{\geq0}$. Define the decreasing filtration
\begin{equation}
\mathscr{F}^jC_\ast^\mathit{nm}(k):=\bigcap_{0\leq i\leq j}\ker(\sigma_{k-1,i}),\textrm{ }0\leq j\leq k-1,
\end{equation}
and the increasing filtration
\begin{equation}
\mathscr{F}^{j+1}D_\ast(k):=\mathrm{span}_\mathbb{R}\left\langle\mathrm{im}(\delta_{k,1}),\cdots,\mathrm{im}(\delta_{k,j+1})\right\rangle,\textrm{ }0\leq j\leq k-1.
\end{equation}
We will argue by induction on $j$. First note that for $j=0$ the short exact sequence
\begin{equation}
0\rightarrow\mathscr{F}^1D_\ast(k)\rightarrow C_\ast(k)\rightarrow\mathscr{F}^0C_\ast^\mathit{nm}(k)\rightarrow0
\end{equation}
splits by the cosimplicial identity $\sigma_{k-1,0}\circ\delta_{k,1}=\mathit{id}_C$, where the map from $\mathscr{F}^1D_\ast(k)$ to $C_\ast(k)$ is the obvious inclusion, and the map from $C_\ast(k)$ to $\mathscr{F}^0C_\ast^\mathit{nm}(k)$ is defined by $x\mapsto x-\delta_{k,1}\circ\sigma_{k-1,0}(x)$.

Assume that we already have an isomorphism $\mathscr{F}^iC_\ast^\mathit{nm}(k)\cong\frac{C_\ast(k)}{\mathscr{F}^{i+1}D_\ast(k)}$ for $0\leq i\leq j-1$ and all $k\in\mathbb{N}$, there is a commutative diagram
\begin{equation}\label{eq:diagram}
\begin{tikzcd}
	0 \arrow[r] &\mathscr{F}^{j-1}C_\ast^\mathit{nm}(k-1) \arrow[d,"\cong"] \arrow[r,"\delta_{k,j+1}"] &\mathscr{F}^{j-1}C_\ast^\mathit{nm}(k) \arrow[d, "\cong"] \arrow[r,"p_{k,j}"] &\mathscr{F}^jC_\ast^\mathit{nm}(k) \arrow[d] \arrow[r] &0 \\
	0 \arrow[r] &\frac{C_\ast(k-1)}{\mathscr{F}^{j}D_\ast(k-1)} \arrow[r, "\delta_{k,j+1}"] &\frac{C_\ast(k)}{\mathscr{F}^{j}D_\ast(k)} \arrow[r,"\mathit{pr}"] &\frac{C_\ast(k)}{\mathscr{F}^{j+1}D_\ast(k)} \arrow[r] &0
\end{tikzcd}
\end{equation}
where the map $p_{k,j}$ is given by $x\mapsto x-\delta_{k,j+1}\circ\sigma_{k-1,j}(x)$ and $\mathit{pr}$ is the standard projection. 

To see this, we first check that the upper row is exact. By (\ref{eq:sd}), $\sigma_{k-1,i}\circ\delta_{k,j+1}(x)=\delta_{k-1,j}\circ\sigma_{k-2,i}(x)=0$ for $0\leq i\leq j-1$ and $x\in\mathscr{F}^{j-1}C_\ast^\mathit{nm}(k-1)$. Moreover, $\sigma_{k-1,i}(x)-\sigma_{k-1,i}\circ\delta_{k,j+1}\circ\sigma_{k-1,j}(x)=0$ for $0\leq i\leq j$ and $x\in\mathscr{F}^{j-1}C_\ast^\mathit{nm}(k)$ by the cosimplicial identities (\ref{eq:cos}) and (\ref{eq:sd}). Thus the upper row is well-defined. Its exactness follows from
\begin{equation}
\delta_{k,j+1}(x)-\delta_{k,j+1}\circ\sigma_{k-1,j}\circ\delta_{k,j+1}(x)=\delta_{k,j+1}(x)-\delta_{k,j+1}(x)=0,
\end{equation}
where $x\in\mathscr{F}^{j-1}C_\ast^\mathit{nm}(k-1)$ and the first equality follows from (\ref{eq:sd}). The exactness of the lower row is automatic once we show the map $\delta_{k,j+1}$ is well-defined and injective. To see this, let $x=\delta_{k-1,i}(y)$ for some $1\leq i\leq j$ and $y\in\mathscr{F}^{j}D_\ast(k-2)$, then by the cosimplicial identity (\ref{eq:cod}) we have $\delta_{k,j+1}\circ\delta_{k-1,i}(y)=\delta_{k,i}\circ\delta_{k-1,j}(y)\in\mathscr{F}^jD_\ast(k)$. For the injectivity, note that if $\delta_{k,j+1}(x)\in\mathscr{F}^jD_\ast(k)$ for some $x\in C_\ast(k-1)$, then $x=\sigma_{k-1,j}\circ\delta_{k,j+1}(x)\in\mathscr{F}^{j}D_\ast(k-1)$ by (\ref{eq:sd}). Now the first square of the diagram (\ref{eq:diagram}) commutes since it consists of the natural inclusions and projections. For the second square, note that the extra term $\delta_{k,j+1}\circ\sigma_{k-1,j}$ in the definition of $p_{k,j}$ does not affect things modulo $\mathscr{F}^jD_\ast(k)$, so it commutes as well.

Since the horizontal lines of (\ref{eq:diagram}) are short exact sequences, and the first two vertical arrows are isomorphisms, so is the third vertical arrow.
\end{proof}

\end{document}
